\chapter{Hoe computers werken}

Wist u dat veel computerexperts en programmeurs al leerden omgaan met computers toen ze nog kleine kinderen waren? Thuiscomputers werden pas in het begin van de jaren tachtig gemeengoed. Ze waren toen zo nieuw dat mensen vaak zelf programma's schreven om te doen wat zij wilden doen, omdat er nog geen software bestond die dat werk voor hen deed.

Het was ook heel gewoon dat mensen met allerlei kantoorbanen leerden programmeren op de computers die zij voor hun werk gebruikten. Zo leerden bijvoorbeeld de mensen die de salarisadministratie van een bedrijf deden, de computer te programmeren zodat die het salaris van iedereen berekende!

Sindsdien is er veel veranderd. De meeste mensen kiezen tegenwoordig bestaande programma's of toepassingen om te doen wat zij nodig hebben en denken dat programmeren een gespecialiseerde vaardigheid is die slechts sommige mensen kunnen leren. Dat is echter niet waar. Natuurlijk is iedereen, net als in elk ander vakgebied, in sommige dingen beter dan in andere, of het nu om sport, breien, wiskunde of schrijven gaat. Bijna iedereen is in staat genoeg te leren om zich in het leven te kunnen redden.

Wij hebben de MEGA65 gemaakt, omdat wij ervan overtuigd zijn dat juist u kunt leren programmeren, zodat computers nuttiger voor u worden, en dat u, net als bij het aanleren van elke nieuwe vaardigheid, de voldoening, het plezier en de nieuwe avonturen kunt ervaren die dat met zich meebrengt!


\phantomsection
\section{Computers zijn dom. Echt dom}

Hoe kan dit nu waar zijn? Computers kunnen zoveel verschillende dingen doen, vaak duizenden malen sneller dan een mens dat kan. Hoe kunnen wij dan beweren dat computers dom zijn? Het antwoord is dat geen enkele computer iets kan doen wat een persoon hem niet heeft opgedragen. Zelfs aan de nieuwste kunstmatige-intelligentiesystemen heeft men geleerd hoe zij moeten leren (of hoe zij moeten leren, hoe zij moeten leren). Om te begrijpen waarom dit zo is, helpt het om te begrijpen hoe computers werkelijk werken.

\subsection{Een eierdopjescomputer maken}

Het hart van een computer is de Central Processing Unit (CPU). Veel moderne computers hebben meer dan één processor, maar laten we het om te beginnen eenvoudig houden. De processor heeft een reeks eenvoudige instructies die hij begrijpt, zoals ``haal het ding uit eierdopje \#21'', ``stop dit ding in eierdopje \#403'', ``tel deze dingen bij elkaar op'' of ``voer de volgende instructie uit, maar alleen als het ding in eierdopje \#712 het getal 3 is.''

Wat bedoelen we eigenlijk met al die ``dingen'' en ``eierdopjes''? Laten we beginnen met na te denken over hoe we ons als computer kunnen voordoen, met niets meer dan een lege eierdoos, wat kleine papiertjes en een potlood of pen. Schrijf eerst in elk van de kleine eierdopjes van de eierdoos een getal, te beginnen bij één. Schrijf daarna het getal nul op een klein papiertje en leg dat papiertje in het eerste eierdopje. Doe hetzelfde voor de andere eierdopjes. U heeft nu een eierdoos met genummerde eierdopjes, waarin elk eierdopje een papiertje bevat waarop het getal nul staat. Nu hoeven we alleen nog een paar eenvoudige regels te bedenken die beschrijven hoe onze eierdopjescomputer werkt:

\begin{itemize}
  \item Ten eerste mag elk eierdopje precies één papiertje tegelijk bevatten. Nooit meer. Nooit minder. Zo weten we zeker dat er, wanneer we de vraag stellen ``wat zit er in eierdopje zoveel'', één eenduidig antwoord is. Zo werkt ook het computergeheugen: elke geheugenlocatie kan slechts één ding tegelijk bevatten.

  \item Ten tweede heeft de computer een manier nodig om te weten wat hierna moet gebeuren. Op de meeste computers heet dit de programmateller (PC), niet te verwarren met de PC als afkorting van {\bf p}ersonal {\bf c}omputer. De programmateller is niets meer dan het nummer van de volgende geheugenlocatie (of in ons geval, het eierdopje) die de computer bekijkt wanneer hij beslist wat hierna moet gebeuren. Het is handig om nog een extra papiertje bij de hand te hebben waarop u onderweg het nummer van de PC kunt schrijven.

  \item Ten derde hebben we een lijst nodig van de dingen die de eierdopjescomputer zal doen, op basis van het getal dat in het door de PC aangegeven eierdopje staat.
\end{itemize}

Laten we nu bedenken wat de computer allemaal kan doen, op basis van het getal in het door de PC aangegeven eierdopje. We houden het eenvoudig en beperken ons tot het volgende:

\begin{center}
  \begin{longtable}{|R{2.2cm}|p{8cm}|}

    \hline
        {\textbf{Getal in het eierdopje}} & {\textbf{Handeling}} \\ \hhline{|=|=|}        0 & {i) Tel één op bij de PC en doe verder niets.} \\ \hline        1 & {i) Tel één op bij de PC. \hfill\break ii) Stel de PC in op het getal dat in dat eierdopje staat.} \\ \hline        2 & {i) Tel één op bij de PC. \hfill\break ii) Tel het getal in het door de PC aangegeven eierdopje op bij het getal in het eierdopje dat wordt aangegeven door het getal in het daaropvolgende eierdopje. \hfill\break iii) Zet het antwoord in het eierdopje dat wordt aangegeven door het daaropvolgende eierdopje. \hfill\break iv) Tel ten slotte nog twee op bij de PC, om de eierdopjes die we gebruikt hebben over te slaan.}  \\ \hline
  \end{longtable}
\end{center}

Maakt u zich geen zorgen als dit nu wat verwarrend klinkt, vooral die laatste -- die nemen we heel binnenkort uitgebreid door! De beste manier om dit uit te leggen is enkele voorbeelden door te werken.

